% Tree (tomato plant) experiment — experimental setup.
% Written to match the style of simulation_environment.tex; demote \section
% to \subsection if it lives under a larger Experiments chapter.

\section{Plant Experiment: Targeted Perception under Natural Occlusion}
\label{sec:tree-setup}

While the panel experiments manufacture occlusion externally, plants occlude
themselves: leaves and branches naturally hide the fruit. To evaluate the
planners under such \emph{natural} occlusion, the target object is replaced
by a single 3D tomato plant mesh (\texttt{tomato6.dae}) composed of ten named
submeshes -- one branch structure, two leaf groups, three blossoms, and four
fruits. The plant is placed at $(0.50,\ {-0.50},\ 0.90)$\,m at a uniform scale
of $0.4$, giving a canopy that spans \SIrange{0.90}{1.49}{\metre} in height
($\approx\SI{0.59}{\metre}$ tall) with a lateral radius of at most
\SI{0.17}{\metre}. The individual fruits measure between \SI{20}{\milli\metre}
and \SI{74}{\milli\metre} across. No artificial occluders are present; the
occluder manifest is empty and all occlusion arises from the plant's own
foliage.\\

\paragraph{Attention levels.} Following the region-of-interest levels of
\citet{burusa2024} (whole plant, main stem, leaf nodes), three perception
tasks of increasing selectivity are defined. Each level determines (i) which
submeshes constitute the reconstruction target and hence the ground truth,
(ii) the ROI cube, and (iii) the reconstruction volume, since the whole plant
does not fit the default grid. Within a level, both planners share identical
settings:

\begin{table}[h]
\centering
\begin{tabular}{lllll}
\toprule
Level & Target submeshes & ROI half-edge & Grid (m) & Voxel $\rho$ \\
\midrule
whole plant & all ten          & \SI{0.31}{\metre}  & $0.38 \times 0.6 \times 0.62$ & \SI{5}{\milli\metre} \\
blossoms    & 3 blossoms       & \SI{0.22}{\metre}  & $0.30 \times 0.6 \times 0.46$ & \SI{4}{\milli\metre} \\
fruits      & 4 fruits         & \SI{0.105}{\metre} & $0.30 \times 0.6 \times 0.30$ & \SI{3}{\milli\metre} \\
\bottomrule
\end{tabular}
\caption{Attention levels of the plant experiment. The ROI is centred on the
target-class centroid: $(0.50, -0.50, 1.20)$\,m for the whole plant,
$(0.50, -0.50, 1.30)$\,m for the blossoms, and $(0.50, -0.50, 1.16)$\,m for
the fruits.}
\label{tab:tree-levels}
\end{table}

The perception pipeline is unchanged: the target-class submeshes are rendered
in red and detected by the HSV-based semantic segmenter, while non-target
parts keep natural colours (brown branch, green leaves) and are mapped as
occupied but non-target space.\\

\paragraph{Plant orientations.} With a single plant model, scene variation is
obtained by rotating the plant about its own vertical axis, in the spirit of
the twelve orientations per plant used by \citet{burusa2024}. Four yaw
angles $\psi \in \{0^\circ, 90^\circ, 180^\circ, 270^\circ\}$ are evaluated;
each orientation hides a different subset of fruits behind different leaves
and therefore constitutes a distinct natural-occlusion pattern. The
ground-truth point set is rotated identically to the spawned model
(\texttt{TREE\_YAW}). Because the fruit and blossom centroids lie on the
rotation axis, the target point, ROI, camera start pose, and workspace remain
identical across orientations -- only the scene rotates. Results are reported
as mean $\pm$ standard deviation over the four orientations.\\

\paragraph{F1 matching threshold.} The panel experiments use a matching
threshold of $4\rho = \SI{12}{\milli\metre}$. On fruit-scale targets this
saturates the metric: a \SI{12}{\milli\metre} tolerance around the sparse
first-view reconstruction already ``recalls'' nearly every ground-truth point
of a \SIrange{20}{74}{\milli\metre} fruit ($F_1 \approx 1$ from the first
viewpoint). A diagnostic run measured the actual distance between
occupancy-thresholded voxel centres and the mesh surface as
\SI{1.5}{\milli\metre} median and \SI{4.8}{\milli\metre} maximum -- far below
the bunny-calibrated offset. The threshold is therefore tightened to
$2\rho$ per level (\SI{6}{\milli\metre}, \SI{8}{\milli\metre},
\SI{10}{\milli\metre} for the fruit, blossom, and whole-plant levels
respectively; \texttt{F1\_THRESH}), which leaves precision unaffected (the
measured maximum surface offset lies below it) while restoring the
discriminative power of recall.\\

\paragraph{Reachability constraint.} The viewpoint workspace of the panel
experiments is retained, but the blossom and whole-plant ROIs sit
\SIrange{0.14}{0.30}{\metre} higher than the panel target, which places part
of the sampling volume beyond the manipulator's reachable workspace: an
inverse-kinematics reach map of the camera-in-hand configuration yielded no
feasible look-at pose beyond $\approx\SI{1.0}{\metre}$ radial distance from
the UR5e base, while every commanded pose that failed motion planning in
preliminary runs lay at \SI{1.10}{\metre} or farther. A local planner is
insensitive to this: GradientNBV's gradient steps keep the camera near
previously executed -- hence feasible -- poses. A globally sampling planner,
by contrast, is systematically \emph{attracted} to the infeasible region:
viewpoints that can never be executed are never observed, so the expected
information gain of the volume they cover remains permanently high, and each
failed motion consumes one viewpoint of the budget while the camera stays in
place. RH-NBV therefore projects every candidate whose distance from the
manipulator base exceeds $R = \SI{1.0}{\metre}$ radially back onto that
reachable shell before evaluation. The constraint encodes the physical reach
of the platform rather than a tuning choice; the information-gain objective
and all remaining parameters are untouched, and GradientNBV, whose poses
remain within reach, is unaffected.\\

\paragraph{Protocol.} Each (level, orientation) scenario is executed on a
freshly restarted simulation stack. Both planners observe the scene for a
budget of $20$ viewpoints from the same initial pose, one trial per
orientation. The RH-NBV hyper-parameters are kept exactly as in the panel
experiments ($K = 10$ candidates, horizon $H = 3$, $\gamma = 0.85$,
$\lambda = 2.0$); no per-object retuning is performed, so any performance
difference reflects the planner's ability to generalise to a new object
class. Reported metrics are ROI coverage, precision, recall and $F_1$ against
the target-class ground truth, the number of views required to reach fixed
coverage thresholds, and the total end-effector trajectory length.
